\documentclass[a4paper]{scrartcl}

% i18n
\usepackage[utf8]{inputenc}
\usepackage[T1]{fontenc}
\usepackage{lmodern}

% layout
\usepackage[cm]{fullpage}
\parindent0mm

% math
\usepackage{mathtools}
\usepackage{amsthm}
\usepackage{amssymb}

% references
\usepackage{hyperref}


% environment for "exercise"
\newcounter{exc}
\newenvironment{exercise}[1]%
{\refstepcounter{exc}\textbf{Exercise 1.\arabic{exc}} \emph{#1}\\}
{

\hrulefill\medskip}%
\renewcommand{\labelenumi}{(\alph{enumi})}

\theoremstyle{definition}
\newtheorem{definition}{Definition}

\newcommand{\N}{\mathbb{N}}
\newcommand{\R}{\mathbb{R}}
\newcommand{\Z}{\mathbb{Z}}
\newcommand{\Q}{\mathbb{Q}}
\newcommand{\vc}[1]{\boldsymbol{#1}}

\title{Exercise Sheet 1: TSP}
\subtitle{Mathematical Programming and Optimization in Transport Logistics}
\author{Student 1 (student id) \and Student 2 (student id) \and Student 3 (student id)}
\date{SS 2026}

\begin{document}

\maketitle

\hrulefill\medskip

\begin{itemize}
	\item Typeset your solution (in, e.g., \LaTeX). You may use the provided \texttt{.tex} template for this.
	\item Cite all references you used, including any use of generative AI.
	\item If you work as a team, list all team members as coauthors in the PDF and the source code (e.g., \texttt{pyproject.toml})
	\item Upload your solution in TUWEL before the deadline. It should consist of \textbf{two files}
		\begin{itemize}
			\item a \texttt{.pdf} with your solution to the exercises, and
			\item an archive (e.g., \texttt{.zip}, \texttt{.tar.gz}) containing your implementation.
		\end{itemize}
	\item Upload \textbf{one solution} per team.
\end{itemize}

\hrulefill\medskip

\begin{exercise}{TSP (\textbf{5 points})}
	\begin{definition}[Travelling Salesman Problem, TSP]
		Given
		a set of cities $N = \{1, \dots, n\}$
		and costs $c_{ij} \geq 0,\ \forall i,j \in N,\ i \neq j,$ for travelling from city $i$ to city $j$,
		find a minimum-cost tour that visits every city exactly once and returns to the starting city at the end.
	\end{definition}
	
	
	Consider the following three formulations for the TSP (as discussed in lecture \texttt{02\_modelling.pdf}).
	W.l.o.g, we choose city $1$ as the start of our tour.
	
	
	\paragraph{Sequential formulation (SEQ)}
	
	Binary variables $x_{ij} \in \{0, 1\}, \forall i,j \in N,\ i \neq j$, indicate whether we travel directly from city $i$ to city $j$.
	Continuous variables $u_i \in [1, n-1], \forall i \in N \setminus \{1\}$, indicate the position of city $i$ within the selected tour (Note that city $1$ is always implicitly at position $0$).
	
	\begin{align}
		\min\ & \sum_{i,j \in N: i \neq j} c_{ij} x_{ij} \label{eq:seq:obj} \\
		% common constraints
		\text{s.t.}\ & \sum_{j \in N \setminus \{i\}} x_{ji} = 1 \label{eq:seq:incoming-arc} & \forall i \in N \\
		& \sum_{j \in N \setminus \{i\}} x_{ij} = 1 & \forall i \in N \label{eq:seq:outgoing-arc} \\
		% SEQ constraints
		& u_i + x_{ij} \leq u_j + (n-2) \cdot (1 - x_{ij}) & \forall i,j \in N \setminus \{1\},\ i \neq j \label{eq:seq:sequencing} \\
		% variable domains
		& x_{ij} \in \{0, 1\} & \forall i,j \in N,\ i \neq j \label{eq:seq:x-domain} \\
		& 1 \leq u_i \leq n - 1 & \forall i \in N \setminus \{1\} \label{eq:seq:u-domain}
	\end{align}
	
	The objective function \eqref{eq:seq:obj} minimizes the total cost of all selected travel legs.
	Constraints \eqref{eq:seq:incoming-arc} and \eqref{eq:seq:outgoing-arc} ensure that each city is entered and exited exactly once, respectively.
	
	Constraints \eqref{eq:seq:sequencing} ensure that if we travel from city $i$ to city $j$, city $j$'s position in the tour is \emph{later} than city $i$'s.
	If we do \emph{not} travel from city $i$ to city $j$, the corresponding constraint is deactivated by the second term on the RHS (which ensures that the constraint is always trivially satisfied, irrespective of how variables $u_i$ and $u_j$ are set).
	These constraints prevent subtours by disallowing any cycles not containing node $1$, since the cities in these cycles would have to be assigned ever-increasing tour position variables $u_i$.
	
	Constrains \eqref{eq:seq:x-domain} and \eqref{eq:seq:u-domain} define the domains of variables $x_{ij}$ and $u_i$, respectively.
	
	
	\paragraph{Single-commodity flow formulation (SCF)}
	
	Binary variables $x_{ij} \in \{0, 1\}, \forall i,j \in N,\ i \neq j$, indicate whether we travel directly from city $i$ to city $j$.
	Continuous variables $f_{ij} \in [0, n-1], \forall i,j \in N,\ i \neq j$, indicate the amount of ``flow'' from city $i$ to city $j$.
	
	\begin{align}
		\min\ & \sum_{i,j \in N: i \neq j} c_{ij} x_{ij} \label{eq:scf:obj} \\
		% common constraints
		\text{s.t.}\ & \sum_{j \in N \setminus \{i\}} x_{ji} = 1 \label{eq:scf:incoming-arc} & \forall i \in N \\
		& \sum_{j \in N \setminus \{i\}} x_{ij} = 1 & \forall i \in N \label{eq:scf:outgoing-arc} \\
		% SCF constraints
		& \sum_{j \in N \setminus \{1\}} f_{1j} = n - 1 \label{eq:scf:flow-generation} \\
		& \sum_{j \in N \setminus \{i\}} f_{ji} - \sum_{j \in N \setminus \{i\}} f_{ij} = 1 & \forall i \in N \setminus \{1\} \label{eq:scf:flow-consumption} \\
		& f_{ij} \leq (n-1) \cdot x_{ij} & \forall i,j \in N,\ i \neq j \label{eq:scf:linking} \\
		% variable domains
		& x_{ij} \in \{0, 1\} & \forall i,j \in N,\ i \neq j \label{eq:scf:x-domain} \\
		& 0 \leq f_{ij} \leq n - 1 & \forall i,j \in N,\ i \neq j \label{eq:scf:f-domain}
	\end{align}
	
	The objective function \eqref{eq:scf:obj} minimizes the total cost of all selected travel legs.
	Constraints \eqref{eq:scf:incoming-arc} and \eqref{eq:scf:outgoing-arc} ensure that each city is entered and exited exactly once, respectively.
	
	Constraint \eqref{eq:scf:flow-generation} ensures that exactly $n-1$ units of flow leave city $1$,
	while constraints \eqref{eq:scf:flow-consumption} ensure that each other city consumes exactly one unit of flow.
	Constraints \eqref{eq:scf:linking} ensure that flow may only be sent between two cities $i$ and $j$ if we travel directly between them.
	Together, these constraints prevent subtours by forcing all cities to be connected to city $1$, since the flow consumed by each city must originate at city $1$ and be routed to city $i$ only via city-to-city connections selected by the corresponding $x_{ij}$ variables.
	
	Constrains \eqref{eq:scf:x-domain} and \eqref{eq:scf:f-domain} define the domains of variables $x_{ij}$ and $f_{ij}$, respectively.
	
	
	\paragraph{Multi-commodity flow formulation (MCF)}
	
	Binary variables $x_{ij} \in \{0, 1\}, \forall i,j \in N,\ i \neq j$, indicate whether we travel directly from city $i$ to city $j$.
	Continuous variables $f_{ij}^k, \forall i,j \in N,\ i \neq j$, indicate the amount of ``flow'' of type $k$ from city $i$ to city $j$. Intuitively, flow of type $k$ is meant to go from city $1$ to city $k$.
	
	\begin{align}
		\min\ & \sum_{i,j \in N: i \neq j} c_{ij} x_{ij} \label{eq:mcf:obj} \\
		% common constraints
		\text{s.t.}\ & \sum_{j \in N \setminus \{i\}} x_{ji} = 1 \label{eq:mcf:incoming-arc} & \forall i \in N \\
		& \sum_{j \in N \setminus \{i\}} x_{ij} = 1 & \forall i \in N \label{eq:mcf:outgoing-arc} \\
		% MCF constraints
		& \sum_{j \in N \setminus \{1\}} f_{1j}^k - \sum_{j \in N \setminus \{1\}} f_{j1}^k = 1 & \forall k \in N \setminus \{1\} \label{eq:mcf:flow-generation} \\
		& \sum_{j \in N \setminus \{k\}} f_{kj}^k - \sum_{j \in N \setminus \{k\}} f_{jk}^k = -1 & \forall k \in N \setminus \{1\} \label{eq:mcf:flow-consumption} \\
		& \sum_{j \in N \setminus \{i\}} f_{ij}^k - \sum_{j \in N \setminus \{i\}} f_{ji}^k = 0 & \forall i,k \in N \setminus \{1\},\ i \neq k \label{eq:mcf:flow-conservation} \\
		& f_{ij}^k \leq x_{ij} & \forall i,j \in N,\ i \neq j,\ k \in N \setminus \{1\} \label{eq:mcf:linking} \\
		% variable domains
		& x_{ij} \in \{0, 1\} & \forall i,j \in N,\ i \neq j \label{eq:mcf:x-domain} \\
		& 0 \leq f_{ij}^k \leq 1 & \forall i,j \in N,\ i \neq j,\ k \in N \setminus \{1\} \label{eq:mcf:f-domain}
	\end{align}
	
	The objective function \eqref{eq:mcf:obj} minimizes the total cost of all selected travel legs.
	Constraints \eqref{eq:mcf:incoming-arc} and \eqref{eq:mcf:outgoing-arc} ensure that each city is entered and exited exactly once, respectively.
	
	Constraints \eqref{eq:mcf:flow-generation} ensure that city $1$ generates one unit of each type of flow $k$, whereas constraints \eqref{eq:mcf:flow-consumption} ensure that city $k$ consumes exactly that unit of $k$-flow.
	Any flow that enters some city $i$ that is not meant to be consumed there (i.e., if $i \neq k$) is simply sent onward (due to constraints \eqref{eq:mcf:flow-conservation}).
	Constraints \eqref{eq:mcf:linking} ensure that flow may only be sent between two cities if we travel directly between them.
	This prevents subtours in much the same way as the SCF formulation by ensuring that all cities are connected to city $1$ (since the flow each city consumes must originate there and can only be sent between two cities if we travel between them).
	
	Constrains \eqref{eq:mcf:x-domain} and \eqref{eq:mcf:f-domain} define the domains of variables $x_{ij}$ and $f_{ij}^k$, respectively.
	
	
	\bigskip
	\begin{itemize}
		\item \textbf{Implement} each formulation (SEQ, SCF, MCF) with an ILP modelling framework.
		
		We recommend you use the provided Python template (using \texttt{gurobipy} as its modelling framework and Gurobi as its ILP solver -- see TUWEL for how to get a license), but you may use any (reasonable) programming language and ILP modelling framework / solver you want.
		Please not, however, that we do not provide any support for other programming languages, modelling frameworks or ILP solvers.
		
		\item Use each implementation to \textbf{solve} the three provided instances
		(\texttt{ulysses16.tsp}, \texttt{dantzig42.tsp}, and \texttt{att48.tsp}).
		For each instance and formulation, state
		\begin{itemize}
			\item the number of variables (total, continuous, integer, binary),
			\item the number of constraints,
			\item the required runtime and maximum memory consumption,
			\item the number of branch-and-bound nodes explored,
			\item the optimal objective value, and
			\item an optimal tour (i.e., the sequence in which the cities are visited).
		\end{itemize}
		
		Note that solving some instances with some models may take a few minutes.
	\end{itemize}
\end{exercise}


\end{document}
